%%
% Author Thomas C. Hales
% LaTeX Format


%!TEX TS-program = latex    
%% This line is for TexShop. 

% Revision history. See svn.
% Versions started Jan 2007.

% Revisions Sept 2009: dependency on Tarski removed. 2 enclosed over quad removed. hypermap algorithm rewritten.  

% 2012-1-25, svn 2659, sent to Astley CUP by gmail.
% "Final Version"

% 2012-5-3, svn 2787, CUP copy, gmail and printed copies sent.
% version 2787 is the post-galley-proof version.

% 2012-5-29 svn 2845, CUP copy, thunderbird copy sent.
% 


%% XX unable to make the \part command work.
%% In pdf sidebar table of contents get stray markings in Appendix.
%% XX unable to get endnotes to work

\documentclass[spanningrule]{cambridge7A}
%% Cambridge University Press Macros from
%% https://authornet.cambridge.org/information/productionguide/laTex_files/


% required by CUP
\usepackage[numbers]{natbib}
\usepackage{rotating}
\usepackage{floatpag}
 \rotfloatpagestyle{empty}
\usepackage{amsthm}
\usepackage{graphicx}
\usepackage{multind}\ProvidesPackage{multind}
\usepackage{times}

% my additions
\usepackage{verbatim}
\usepackage{latexsym}
\usepackage{amsfonts}
\usepackage{amsmath}
\usepackage{crop}
\usepackage{txfonts}
\usepackage[hyphens]{url}
\usepackage{setspace}
\usepackage{ellipsis} % 
% http://www.ctan.org/tex-archive/macros/latex/contrib/ellipsis/ellipsis.pdf 
% wrapping graphics
% http://en.wikibooks.org/wiki/LaTeX/Floats,_Figures_and_Captions
\usepackage{wrapfig}


%tikz graphics
\usepackage{xcolor} % to remove color.
\usepackage{tikz} % Needs pgf version >= 2.10.
\usetikzlibrary{chains,shapes,arrows,shadings,trees,matrix,positioning,intersections,decorations,
  decorations.markings,backgrounds,fit,calc,fadings,decorations.pathreplacing}
% fonts
\usepackage[mathscr]{euscript} % powerset.
\usepackage{pifont} %ding
\usepackage[displaymath]{lineno}


%%%%%%%%%%%%%%%%%%%%%%%%%%%%%%%%%%%%%%%%%

% This file contains local settings and system dependencies



% Auxiliary directories
\def\ps{../../../../Pictures/kep07ps}  % flypaper graphics
\def\pdf{/Users/thomashales/Pictures/collect_geom} % tarski graphics

\def\showgraphics{f}  
% t: display graphics (there are none to show yet)
% f (default): print a "no graphics logo" where graphics would normally go.


\def\displayallproof{t} 
% t (default): display all proofs.
% f: print documents without the proofs-- theorem statements only

\def\displayrating{f}
% t (default): display all ratings (verbose is also true)
% f : don't show them.

\def\verbose{f}
% f (default): do not display debugging information,
% t : display debug information and information about the formalization.

\def\tarskipagesep{f}
% f (default): no page breaks between lemmas in tarski collection.
% t :add page breaks

\def\tarskistrip{f}
% f (default): preserves control sequences in tarski collection
% t: control sequences in tarski collection are stripped and printed to an auxiliary file.

\def\tarskidump{f}
% t : dump tarski guid data to external write file.
% f (default) : don't use external write file.
 
%-%
% --Repository--
%-%
% generate revision number by
% svn propset svn:keywords "LastChangedRevision" kepmacros.tex
\def\svninfo{%
  TeXed on \today; \hfill\break
  Repository Root: https://flyspeck.googlecode.com/svn \hfill\break
  SVN $LastChangedRevision$
  }

%-%
% --Fonts--
%-%
\font\twrm=cmr8

%-%
% --Graphics--
%-%
%set \showgraphics option in flag_fly.tex
% flypaper graphics
\def\myincludegraphics#1{%
      \if\showgraphics t{\includegraphics{#1}}%
      \else{\includegraphics{noimage.eps}}\fi}
\def\szincludegraphics[#1]#2{%
      \if\showgraphics t{\includegraphics{#2}}%
      \else{\includegraphics{noimage.eps}}\fi}

% kepler graphics
\def\pdffigtemplatex#1#2#3{%
\begin{figure}[htb]%
  \centering
  \myincludegraphics{\pdfp/#1}
  \caption{#3}
  \label{fig:#2}%
\end{figure}%
}
\def\pdfg#1#2#3{\if\showgraphics t{\pdffigtemplatex{#1}{#2}{#3}}\else{}\fi}


% tarski graphics
\def\pdffigtemplate#1#2#3{%
\begin{figure}[htb]%
  \centering
  \myincludegraphics{\pdf/#1}
  \caption{#3}
  \label{tarski:fig:#2}%
\end{figure}%
}
\def\pdffig#1#2#3{\if\showgraphics t{\pdffigtemplate{#1}{#2}{#3}}\else{}\fi}

%-%
% --Footnotes--
%-%
% http://help-csli.stanford.edu/tex/latex-footnotes.shtml
\long\def\symbolfootnote[#1]#2{\begingroup%
\def\thefootnote{\ensuremath{\fnsymbol{footnote}}}\footnote[#1]{#2}\endgroup}

%-%
% --Special Formatting--
%-%
% http://en.wikibooks.org/wiki/LaTeX/Formatting#List_Structures
\renewcommand{\labelitemii}{$\star$}


%-%
% --Indexing, References, Citations--
%-%
\def\indy#1#2{\index{index/#1}{#2}}

%-%
% --Proof Display--
%-%
% set with \displayallproof in flag_fly. If f, then proofs are swallowed.
%% "proved" environment. toggle with \displayallproof
%
\def\hide#1{}
\def\swallowed{\relax}
\def\swallow#1\swallowed{}
\newenvironment{iproved}{}{}
\newenvironment{proved}{\resetproved\begin{iproved}}{\end{iproved}}
\def\hideproof{\renewenvironment{iproved}{%
   \centerline{\it -- Proof Proofed --}
  \renewenvironment{itemize}{}{}
  \renewenvironment{enumerate}{}{}
  \def\item{\relax}
  \catcode13=12
  \swallow
}{}}
\def\showproof{\renewenvironment{iproved}{\begin{proof}}{\end{proof}}}
\def\resetproved{\if\displayallproof t\showproof\else\hideproof\fi}



%-%
% --Debugging Information--
%-%
\def\uf#1{{\par\narrower\it #1}} % unfinished manuscript
%% verbose:
\def\rating#1{\if\displayrating t{{\textsc {[rating={\ensuremath {#1}}].\ }}}\else{}\fi}
\def\oldrating#1{\if\displayrating t{{\textsc {[former rating={\ensuremath {#1}}].\ }}}\else{}\fi}
\def\formal#1{\if\verbose t{{\tt [formal: #1].\ }}\else{}\fi}
\def\formalauthor#1{\if\verbose t{{\tt [formal proof by: #1].\ }}\else{}\fi}
\def\footformal#1{\if\verbose t{\footnote{\formal{#1}}}\else{}\fi}
\def\usage#1{\if\verbose t{\symbolfootnote[1]{#1}}\else{}\fi}
\def\move#1#2{\begin{#1} See \ref{#2}.\end{#1}}
\def\FIXX#1{\if\verbose t{\footnote{\tt [#1]}}\else{}\fi}
\def\page#1#2{#2 {\it [dcgp.#1]}}
% first entry Lemma or Sec of DCG, second entry page.
\def\dcg#1#2{{\if\verbose t{{\tt{[DCG-#1]}}\indy{References}{ZC{#2 #1}@{DCG-#1}|page{#2}}}\else{}\fi}}
\def\tlabel#1{\label{#1}\if\verbose t{{\tt [#1].\ }%
   \indy{References}{#1|itt}}\else{}\fi}
\def\ifverbose#1{\if\verbose t{{#1}}\else{}\fi}




%% Indexing
\def\refXX{0} % for unknown links.
\def\showref#1{\ref{#1}{\tt[#1]}\indy{References}{#1}}
\def\oldlabel#1{\label{x-#1}}

\def\tref#1{\ref{#1}\indy{References}{#1}}
\def\itt#1{{\bf #1}}
\def\guid#1{{\tt[#1].\ }\indy{References}{ZA{#1}@{#1}|itt}}
% textsc
\def\calc#1{{\textsc{calc-#1}}\indy{Interval}{{#1}@{#1}}}
\def\newcalc#1{{\textsc{newcalc-#1}}\indy{Interval}{{#1}@{#1 (new XX)}}}
\def\assembly#1{{\textsc{assembly-#1}}\indy{References}{ZD{#1}@{#1}}}
%\def\conseq#1{{\textsc{conseq-#1}}}

%% 


%% Tarski indexing. Write summary data to external file "\mywrite",
% which is generated in tarski.tex
% Need to extend cs character set.  Use this only within a block.
\def\setcat{
\catcode`\-=11
\catcode`\.=11
\catcode`\:=11
\catcode`\0=11
\catcode`\1=11
\catcode`\2=11
\catcode`\3=11
\catcode`\4=11
\catcode`\5=11
\catcode`\6=11
\catcode`\7=11
\catcode`\8=11
\catcode`\9=11
}
\def\tarfe#1#2{{\foote{{\tt[L.E.G:#1]}~\csname #1-sum\endcsname%
  \expandafter\ifx \csname #1-guid\endcsname \relax [XX-BAD-IDENTIFIER]\else{}\fi
  }
  {\csname #1-guid\endcsname{#2}}}}
\def\tarf#1{\tarfe{#1}{}}
\def\tarfE#1{\tarfe{EE}{-\csname #1\endcsname}}
% control sequences will contain [-.0-9:] \catcode`\-=11, etc.
\newtoks\mysummary
\newtoks\myguid
\newtoks\myname
\def\separator{\if\tarskipagesep t{\clearpage}\else%
  \smallskip\hrule height 0.5pt depth 0pt width 60pt\bigskip\fi}
\def\identity#1{#1}
\def\foote#1#2{\footnote{#1~{\it (#2)}}}
\def\marker#1#2{\par\hangafter1\hangindent=1em\noindent{\bf #1:}\quad #2}
\def\name#1{\marker{Name}{#1}\myname={#1}
  \write\myhtml{  "#1", //name}
}
\def\summary#1{\marker{Summary}#1\mysummary={#1}
  % cleanse control sequences for html
  \if\tarskistrip t{
  \def\sqr{sqrt}
  \def\sqrt{sqrt}
  \def\op{\identity}
  \def\epsilon{epsilon}
  \def\Delta{Delta}
  \def\ups{upsilon}
  \def\beta{beta}
  \def\rho{rho}
  \def\CalE{E}
  \def\chi{chi}
  \def\eta{eta}
  \write\myhtml{  "#1", // summary}
  }
  \else\relax\fi
}
\def\guid#1{\if\tarskidump f{{\tt [#1]}}\else{\marker{ID}#1
   \indy{guid}{#1}
   \immediate\write\mywrite{\global\expandafter\def\csname \the\myname-sum \endcsname{\the\mysummary}}
   \immediate\write\mywrite{\global\expandafter\def\csname \the\myname-guid \endcsname{#1}}
   \immediate\write\myhtml{  "#1" // guid}}%
   \fi
}
\def\tag#1{\marker{Tags}#1
   \immediate\write\myhtml{  "#1", // tag}
}
%\def\rating#1{\marker{Rating}#1
%   \write\myhtml{  <rating>#1</rating>}
%}
\newenvironment{tarski}{%
\write\myhtml{new Tarski(}
%\def\name   {\marker{Name}}
%\def\summary{\marker{Summary}}
%\def\tag    {\marker{Tags}}
\def\rating {\marker{Rating}}
\def\hol{\marker{HOL-Light}}
%\def\guid   {\marker{ID}}
\def\tlabel{\label}
}{\separator\immediate\write\myhtml{),}
}
\newenvironment{tarskidata}{\write\myhtml{var tarskis=[}}
{
\immediate\write\myhtml{];}
}
%% tarski indexing
%\def\tarf#1{\footnote{{\tt geom-#1~\ref{#1}}\indy{ZA{#1}@{#1}}}}




%-%
% --Symbols--
%-%
% norm
\def\|{{\hskip0.1em|\hskip-0.15em|\hskip0.1em}}
\def\mid{\ :\ }
\def\norm#1#2{\|#1 - #2\|}
\def\normo#1{{\|#1\|}}
\def\sland{\ \land\ }
% Sam's
\def\myscorept{\text{ \sl pt}}
\def\qrtet{{quasi-regular tetrahedron}}
\def\qrtets{{quasi-regular tetrahedrons}}
\def\tomcite{{}}
\DeclareMathOperator{\myscore}{\sigma}
 \DeclareMathOperator{\gma}{gma}
 \DeclareMathOperator{\score}{\sigma}
 \DeclareMathOperator{\sol}{sol}
 \DeclareMathOperator{\dih}{dih}
\newtheorem{calcf}{Calculation}[subsection]

% mathcal
\def\CalB{{\mathcal B}}
\def\CalC{{\mathcal C}}
\def\CalD{{\mathcal D}}
\def\CalE{{\mathcal E}}
\def\FF{{\mathcal F}}
\def\CalF{{\mathcal F}}
\def\CalQ{{\mathcal Q}}
%\def\CalW{{\mathcal W}}
\def\CalR{{\mathcal R}}
\def\CalS{{\mathcal S}}
\def\CalV{{\mathcal V}}
\def\T{{\mathcal T}}
\def\Q{{\mathcal Q}}


% brackets
\def\leftopen{]}
\def\leftclosed{[}
\def\rightopen{[}
\def\rightclosed{]}

%% obsolete
\def\leftb{[}
\def\rightb{]}

% squiqqly relations
\def\seq{\approx}
\def\sle{\preceq}
\def\sge{\succeq}
\def\slt{\prec}
\def\sgt{\succ}

% mathbb
\def\R{{\mathbb R}}
\def\N{{\mathbb N}}
\newcommand{\ring}[1]{\mathbb{#1}}
\def\A{{\mathbf A}}
\def\Rp{\ring{R}^{3\,\prime}}

% operatorname
\def\op#1{{\operatorname{#1}}}
\def\optt#1{{\operatorname{{\texttt{#1}}}}}

\def\opat{\op{@}}
\def\atn{\op{arctan\ensuremath{_2}}}
\def\azim{\op{azim}}
\def\sol{\operatorname{sol}}
\def\vol{\op{vol}}
\def\quo{\operatorname{quo}}
\def\anc{\operatorname{anc}}
\def\cro{\operatorname{crown}}
%\def\vor{\operatorname{vor}}
%\def\svor{\operatorname{svor}}
\def\sign{\operatorname{sign}}
\def\octavor{\operatorname{octavor}}
\def\dih{\operatorname{dih}}
\def\Adih{\operatorname{Adih}}
\def\arc{\operatorname{arc}}
\def\cosarc{\operatorname{cosarc}}
\def\rad{\operatorname{rad}}
\def\gap{\operatorname{gap}}
\def\sc{{\operatorname{sc}}}
\def\geom{{\operatorname{g}}}
\def\anal{{\operatorname{an}}}
\def\PM{\operatorname{PM}}
\def\bool{\operatorname{bool}}
\def\true{\op{true}}
\def\false{\op{false}}
\def\flat{\operatorname{flat}}
\def\tangle#1{\langle #1\rangle}
\def\ceil#1{\lceil #1\rceil}
\def\floor#1{\lfloor #1\rfloor}
\def\ups{\upsilonup} % Needs txfonts; else use \upsilon
\def\orgn{\varthetaup} % center of packing
%\def\comp#1{\llbracket #1 \rrbracket}
\def\comp#1{[#1]}
\def\Wdart{W_{\text{dart}}}
\def\Wedge{W_{\text{edge}}}
\def\cell{\operatorname{cell}}

\def\SA{A}
\def\SB{B}
\def\SC{C}
\def\SD{D}
\def\SE{E}
\def\del{\partial}
\def\doct{\delta_{oct}}
\def\dtet{\delta_{tet}}
\def\pt{\hbox{\it pt}}
\def\hm{{h_0}} % 1.26
%\def\Vol{\hbox{vol}}
\def\scoregoal{8\,\pt}
\def\maxpi{\pi_{\max}}
\def\tausc{{\tau\!\operatorname{sc}}}
\def\piF{{\pi_F}}
\def\xiG{\xi_\Gamma}
\def\piG{\pi_\Gamma}
\def\xiV{\xi_V}
\def\xik{\xi_\kappa}
\def\xikG{\xi_{\kappa,\Gamma}}
\def\piV{\pi_V}
\def\tauLP{{\tau_{\hbox{\twrm LP}}}}
\def\DLP{\operatorname{D}_{\hbox{\twrm LP}}}
\def\ZLP{\operatorname{Z}_{\hbox{\twrm LP}}}
\def\tlp{\tau_{\hbox{\twrm LP}}}  % 2 args (p,q) tri, quad
\def\sLP{\sigma_{\hbox{\twrm LP}}}  % 2 args (p,q) tri, quad
\def\squander{(4\pi\zeta-8)\,\pt}
\def\trgt{\operatorname{\it{target}}}
\def\tildeF{{\hbox{$\tilde F$}}}
\def\zloop{{z}_{loop}}
\def\dloop{{\delta}_{loop}}
\def\sqr{\sqrt}
\def\tildeV{V^{\CalS}}
%\def\BigD{\Delta}
%\def\bigd{\Delta^0}
\def\FCinner{\op{FC}_{\text{inner}}}
\def\FC{\op{FC}}
\def\FCR{\op{FCR}}
\def\rogFC{\op{rog}^0_{\text{FC}}}
\def\chino{\chi_0}

\def\x#1{\op{\bf c}{#1}}  % marks context \x(p,r).
\def\xs#1{\op{\bf sc}{#1}}  % marks special context \x(p,r,s).
\def\xtazimlt{\x{(3,1: \op{azim}(\op{gap}) <\pi)}}
\def\xtazimge{\x{(3,1: \op{azim}(\op{gap}) \ge \pi)}}
\def\xfgap{\x{(4,1:\op{gap})}}
\def\y#1{\op{\bf ep}{#1}}  % marks edge parameters \y(n,k).
\def\pqr#1{#1} % marks type (p,q,r).
%\def\pqr#1{\op{\bf typ}{#1}} % marks type (p,q,r).

%% HYPERMAP macros:
% avoid e for both hypermap edge and edge {v,w}
\def\e{\varepsilonup}
\def\ocirc{}



%% FORMULATION macros:
\def\lam{\lambda}
\def\Lam{\Lambda}
\def\bl{{\underline{\lam}}}
\def\bm{{\underline{\mu}}}
\def\angle#1#2#3{{\op{angle}(#1,{#2,#3})}}


\def\card{\op{card}}

% accent used in India citation in Overview.
\def\=#1{\accent"16 #1}
\def\i{I}
%\def\={\relax}



%% Enclosed CALE control sequences:
\expandafter\def\csname 2t0-doesnt-pass-through\endcsname{01}
\expandafter\def\csname E1453\endcsname{02}   
\expandafter\def\csname dcg-p89\endcsname{03}   
\expandafter\def\csname enclosed-v\endcsname{04}   
\expandafter\def\csname 277\endcsname{05}   
\expandafter\def\csname 245\endcsname{06}   
\expandafter\def\csname 245bis\endcsname{07}   
\expandafter\def\csname node\endcsname{08}   
\expandafter\def\csname E:part4:1\endcsname{09}   
\expandafter\def\csname E:part4:2\endcsname{10}   
\expandafter\def\csname E:part4:3\endcsname{11}   
\expandafter\def\csname E:part4:4\endcsname{12}   
\expandafter\def\csname E:part4:5\endcsname{13}   
\expandafter\def\csname E:part4:6\endcsname{14}   
\expandafter\def\csname E:part4:7\endcsname{15}   
\expandafter\def\csname E:part4:8\endcsname{16}   
\expandafter\def\csname E:part4:9\endcsname{17}   
\expandafter\def\csname E:part4:10\endcsname{18}   
\expandafter\def\csname dcg-p142\endcsname{19}   
\expandafter\def\csname convex-quad\endcsname{20}   
\expandafter\def\csname last:E\endcsname{21}
\expandafter\def\csname rem2.7\endcsname{22}     
\expandafter\def\csname pass-anchor\endcsname{23}

%% 
 
 
 
 
 
 
 


% line numbers
\def\lll{\resetlinenumber[1]}
\def\linenumberfont{\normalfont\small\sffamily}

\def\tocpart#1{
  \addcontentsline{toc}{part}{\Large{#1}}}
% (or even \LARGE)
  
\crop

\def\linput#1{\lll\input{#1}}


\raggedbottom  % for now.


%%%%%%%%%%%%%%%%%%%%%%%%%%%%%%%%%%%%%%%%

% new theorems

\theoremstyle{plain}
\newtheorem{theorem}[equation]{Theorem}
\newtheorem{theorem*}[equation]{Theorem$^*$}
\newtheorem{lemma}[equation]{Lemma}
\newtheorem{lemma*}[equation]{Lemma$^*$}
\newtheorem{background}[equation]{Background}
\newtheorem{corollary}[equation]{Corollary}
\newtheorem{example}[equation]{Example}
\newtheorem{conjecture}[equation]{Conjecture}
%---
\theoremstyle{definition}
\newtheorem{definition}[equation]{Definition}
%---
\theoremstyle{remark}
\newtheorem{remark}[equation]{Remark}
\newtheorem{notation}[equation]{Notation}
\newtheorem{exer}[equation]{Exercise}

%%%%%%%%%%%%%%%%%%%%%%%%%%%%%%%%%%%%%%%%%

\begin{document}
\title[a blueprint for formal proofs]
    {%
      Appendices to Dense Sphere Packings}
\author{Thomas C. Hales}

%%%

\indy{Index}{partition|see{decomposition}}%
    
      
%%%%%%%%%%%%%%%%%%%%%%%%%%%%%%%%%%%%%%%%%
%%% FRONT
    \frontmatter
    \maketitle
    \setcounter{page}{5}
    \tableofcontents

%%%%%%%%%%%%%%%%%%%%%%%%%%%%%%%%%%%%%%%%%

   %\noindent

\bigskip\noindent
This material is based upon work supported by the National Science
Foundation under
Grants 0503447 and 0804189.

\bigskip\noindent\svninfo 

%\bigskip\noindent
%A pdf of an earlier version of this work appears at\hfill\break 
%\url{http://www.math.pitt.edu/~thales/papers/Flyspeck_Blueprint_2008.pdf}.

\bigskip\noindent
%% XX Thanks to Tran Nam Trung on the hypermap chapter.

%% INDEX OF FORMAL TERMS.

\def\idx#1#2{\indy{Spec}{{\tt #1}; #2}}
\def\idz#1#2#3{\indy{Spec}{{\tt #1}; #2; {\tt \textbackslash \string#3}}}
\idx{pt}{$\pt$}
\idz{atn2~x~~y}{$\atn(x,y)$}{atn}
\idx{sqrt8}{$\sqrt8$}
\idx{sqrt2}{$\sqrt2$}
\idx{t0}{$t_0$}
\idx{two\_t0}{$2t_0$}
\idx{square\_2t0}{$(2t_0)^2$}
\idx{square\_4t0}{$(4t_0)^2$}
\idx{zeta}{$\zeta$}
\idz{doct}{$\doct$}{doct}
\idz{dtet}{$\dtet$}{dtet}
\idx{pi\_rt18}{$\pi/\sqrt{18}$}
\idx{rogers}{$\sqrt2/\zeta$}
\idx{compression\_cut}{$1.41$}
\idz{squander\_target}{$\squander$}{squander}
\idz{xi\_gamma}{$\xiG$}{xiG}
\idz{xiV}{$\xiV$}{xiV}
\idx{xi'\_gamma}{$\xiG'$}
\idz{xi\_kappa}{$\xik$}{xik}
\idz{xi\_kappa\_gamma}{$\xikG$}{xikG}
\idz{pi\_max}{$\maxpi$}{maxpi}
\idx{t4}{$t_4$}
\idx{t5}{$t_5$}
\idx{t6}{$t_6$}
\idx{t7}{$t_7$}
\idx{t8}{$t_8$}
\idx{t9}{$t_9$}
\idx{t10}{$t_{10}$}
\idx{s5}{$s_5$}
\idx{s6}{$s_6$}
\idx{s7}{$s_7$}
\idx{s8}{$s_8$}
\idx{s9}{$s_9$}
\idx{s10}{$s_{10}$}
\idx{eps0}{$\epsilon_0$}
\idx{Z31}{$Z_{3,1}$}
\idx{D31}{$D_{3,1}$}
\idx{Z31}{$Z(\y{(3,1)})$}
\idx{Z32}{$Z(\y{(3,2)})$}
\idx{Z33}{$Z(\y{(3,3)})$}
\idx{Z41}{$Z(\y{(4,1)})$}
\idx{Z42}{$Z(\y{(4,2)})$}
\idx{D31}{$D(\y{(3,1)})$}
\idx{D32}{$D(\y{(3,2)})$}
\idx{D33}{$D(\y{(3,3)})$}
\idx{D41}{$D(\y{(4,1)})$}
\idx{D42}{$D(\y{(4,2)})$}
\idx{D51}{$D(\y{(5,1)})$}
\idz{delta\_x x1 ... }{$\Delta(x_1,\ldots)$}{Delta}
\idx{delta\_x4 x1 ... }{$\Delta_4(x_1,\ldots)$}
\idx{delta\_x6 x1 ... }{$\Delta_6(x_1,\ldots)$}
\idz{ups\_x x1 ...}{$\ups(x_1,\ldots)$}{ups}
\idz{eta\_x x1 x2 x3}{$\eta(\sqrt{x_1},\sqrt{x_2},\sqrt{x_3})$}{eta}
\idz{eta\_y y1 y2 y3}{$\eta(y_1,y_2,y_3)$}{eta}
\idz{rho\_x x1 ... }{$\rho(x_1,\ldots)$}{rho}
\idx{rad2\_y y1 ... }{$(\rad_V\{v_0,v_1,v_2,v_3\})^2,\ y_1 = |v_0-v_1|,\ldots$}
\idz{chi\_x x1 ... }{$\chi(x_1\ldots)$}{chi}
\idz{dih\_x x1 ... }{$\dih(\sqrt{x_1},\ldots)$}{dih}
\idz{dih\_y y1 ... }{$\dih(y_1,\ldots)$}{dih}
\idx{dih2\_x x1 x2 ... }{$\dih(\sqrt{x_2},\ldots)$}
\idx{dih2\_y x1 x2 ... }{$\dih(y_2,\ldots)$}
\idx{dih3\_x x1 x2 ... }{$\dih(\sqrt{x_3},\ldots)$}
\idx{dih3\_y x1 x2 ... }{$\dih(y_3,\ldots)$}
\idz{sol\_x x1 ... }{$\sol(v_0,\ldots),\ x_1=|v_0-v_1|^2,\ldots$}{sol}
\idz{sol\_y y1 ... }{$\sol(v_0,\ldots),\ y_1=|v_0-v_1|,\ldots$}{sol}
\idx{vol\_x x1 ... }{$\sqrt{\Delta(x_1,\ldots)}/12$}
\idz{beta psi theta}{$\beta_\psi\theta$}{beta}
\idz{arclength a b c}{$\arc(a,b,c)$}{arc}
\idx{volR a b c}{$\op{volR}(a,b,c)$}
\idx{solR a b c}{$\op{solR}(a,b,c)$}
\idx{dihR a b c}{$\op{dihR}(a,b,c)$}
\idx{vorR a b c}{$\op{sovoR}(a,b,c,\lambda_{oct})$}
\idx{denR a b c}{$\delta(a,b,c)$}
\idx{tauR a b c}{$\op{sovoR}(a,b,c,\lambda_{sq})$}
\idx{qy y1 y2 y3 t}{$\op{quovol}(a,b,c)$}
\idx{quo\_x x y z}{}
\idx{qn y1 y2 z t}{}
\idx{phi h t}{}
\idx{phi0}{}
\idx{eta0 h}{}
\idx{crown h}{}
\idx{anc y1 y2 y6}{}
\idx{K0 y1 y2 y6}{}
\idx{AH h t}{}
\idx{BHY y}{}
\idx{KY y1 y2 y3 y4 y5 y6}{}
\idx{KX x1 x2 x3 x4 x5 x6}{}
\idx{vor\_analytic\_x x1 x2 x3 x4 x5 x6}{}
\idx{vor\_analytic\_x\_flipped x1 x2 x3 x4 x5 x6}{}
\idx{octavor\_analytic\_x x1 x2 x3 x4 x5 x6}{}
\idx{tau\_analytic\_x x1 x2 x3 x4 x5 x6}{}
\idx{kappa y1 y2 y3 y4 y5 y6}{}
\idx{kappa\_dih\_y y1 y2 y3 y5 y6 d}{}
\idx{level\_at h t}{}
\idx{vorstar h1 h2 h3 a1 a2 a3 t}{}
\idx{vort\_y y1 y2 y3 y4 y5 y6 t}{}
\idx{vor\_0\_y y1 y2 y3 y4 y5 y6}{}
\idx{tau\_0\_y y1 y2 y3 y4 y5 y6}{}
\idx{vor\_0\_x x1 x2 x3 x4 x5 x6}{}
\idx{tau\_0\_x x1 x2 x3 x4 x5 x6}{}
\idx{vort\_x x1 x2 x3 x4 x5 x6 t}{}
\idx{tauVt\_x x1 x2 x3 x4 x5 x6 t}{}
\idx{vorA\_x x1 x2 x3 x4 x5 x6}{}
\idx{tauA\_x x1 x2 x3 x4 x5 x6}{}
\idx{vorC0\_x x1 x2 x3 x4 x5 x6}{}
\idx{tauC0\_x x1 x2 x3 x4 x5 x6}{}
\idx{vorC\_x x1 x2 x3 x4 x5 x6}{}
\idx{tauC\_x x1 x2 x3 x4 x5 x6}{}
\idx{v0x x1 x2 x3 x4 x5 x6}{}
\idx{v1x x1 x2 x3 x4 x5 x6}{}
\idx{gamma\_x x1 x2 x3 x4 x5 x6}{}
\idx{tau\_gamma\_x x1 x2 x3 x4 x5 x6}{}
\idx{rad2\_x x1 x2 x3 x4 x5 x6}{}
\idx{sigma\_qrtet\_x x1 x2 x3 x4 x5 x6}{}
\idx{sigma1\_qrtet\_x x1 x2 x3 x4 x5 x6}{}
\idx{tau\_sigma\_x x1 x2 x3 x4 x5 x6}{}
\idx{sigma32\_qrtet\_x x1 x2 x3 x4 x5 x6}{}
\idx{mu\_flat\_x x1 x2 x3 x4 x5 x6}{}
\idx{taumu\_flat\_x x1 x2 x3 x4 x5 x6}{}
\idx{mu\_upright\_x x1 x2 x3 x4 x5 x6}{}
\idx{mu\_flipped\_x x1 x2 x3 x4 x5 x6}{}
\idx{vor\_0\_x\_flipped x1 x2 x3 x4 x5 x6}{}
\idx{octavor0\_x x1 x2 x3 x4 x5 x6}{}
\idx{nu\_x x1 x2 x3 x4 x5 x6}{}
\idx{nu\_gamma\_x x1 x2 x3 x4 x5 x6}{}
\idx{taunu\_x x1 x2 x3 x4 x5 x6}{}
\idx{octa\_x x1 x2 x3 x4 x5 x6}{}
\idx{sigmahat\_x x1 x2 x3 x4 x5 x6}{}
\idx{sigmahat\_x' x1 x2 x3 x4 x5 x6}{}
\idx{sigmahatpi\_x x1 x2 x3 x4 x5 x6}{}
\idx{tauhat\_x x1 x2 x3 x4 x5 x6}{}
\idx{tauhatpi\_x x1 x2 x3 x4 x5 x6}{}
\idx{pi\_prime\_tau k0 k1 k2}{}
\idx{pi\_prime\_sigma k0 k1 k2}{}



  % Formal Spec Index.

\smallskip
\newpage


   \newpage
   %%------------------------------------------------------------
% Author: Thomas C. Hales
% Format: LaTeX
% Book Chapter: Dense Sphere Packings
%------------------------------------------------------------

\chapter*{Preface}


%{
%
%\narrower
%
%{\it ``A personality which has been veiled by a formal method
%  throughout many chapters is suddenly seen face to face in the
%  Preface.'' }
%% Introductory note, p3, Famous Prefaces, Harvard Classics, Vol
%% 39. P.F. Collier & Son, 1910.
%
%}

%\bigskip

%\centerline{\it ``Those who justify themselves do not convince.''
%  --Lao-Tzu}
%% quoted in Watts's essay in "Modern Buddhism" ed. Donald Lopez,
%% page 160


{

\narrower

{\it

  ``I think there's a revolution in mathematics around the corner. I
  think that~\dots{}   %in later times
  people will look back on the fin-de-siecle of the twentieth century
  and say `Then is when it happened' (just like we look back at the
  Greeks for inventing the concept of proof and at the nineteenth
  century for making analysis rigorous). I really believe that. And it
  amazes me that no one seems to notice\dots.

  ``Never before have the Platonic mathematical world and the physical
  world been this similar, this close. Is it strange that I expect
  leakage between these two worlds? That I think the proof strings
  will find their way to the computer memories?\dots 

  ``What I expect is that some kind of computer system will be
  created, a proof checker, that all mathematicians will start using
  to check their work, their proofs, their mathematics. I have no idea
  what shape such a system will~\dots{}  take. But I expect some
  system to come into being that is past some threshold so that it is
  practical enough for real work, and then quite suddenly some kind of
  `phase transition' will occur and everyone will be using that
  system.''

{\hfill\rm--Freek Wiedijk~\cite{FWR}} % http://www.cs.ru.nl/~freek/jordan/index.html

}

}

\newpage

{

\narrower\parindent=0pt
\parskip=0.4\baselineskip

{\it

Alecos: Christos has a problem with the `foundational quest'!

Christos:  Wrong!  I have two problems with your  {\rm {version}} of it!  One, it
didn't fail and, two, it wasn't a tragedy!  Granted, there are some tragic
parts!  But the ending is happy, as in the `Oresteia'!  

Apostolos:  Happy for whom?  Cantor, going insane?  G\"odel starving himself to
death out of paranoia? Hilbert or Russell and their psychotic sons? Or Frege with--

Christos: `The meaning
is in the ending!' you said so yourself!  So, follow the quest for ten more years and
you get a brand-new triumphant finale with the creation of the computer, which is the
quest's real hero!   Your problem is, simply, that you see it as a story of people!

Apostolos: Well, stories do tend to be about people!

Christos:  So, choose the right people!  And show what they really did!  All we we learn
of the great von Neumann is he said `It's over'  when he heard G\"odel!

Alecos: But it was over in a sense, wasn't it?  Pop went Hilbert's `no ignorabimus'!

Christos:  But then came the quest's jeune premier, its parsifal~\dots{}  Alan Turing!
He said `Ok, we can't prove everything! So, let's see what we can prove!' and to define
proof, he invented, in 1936, a theoretical machine which contains all the ideas of the
computer!\dots{}  which, after the war, he and von Neumann, the quest's proudest sons,
brought to full life!

{\hfill\rm--Doxiadis and Papadimitriou, ~~Logicomix~\cite{Logicomix}} % page 303.

}

}




\newpage

{

\narrower\parindent=0pt
\parskip=0.4\baselineskip

{\it

``Ever hear of the Kepler Conjecture?''

``Nope.''

I laid the notebook on the table and flipped through the pages. ``It was first stated
in 1611 by Johannes Kepler,'' I said.  ``Kepler became interested in the problem while
he was corresponding with an Englishman named Thomas Harriot, who was trying
to help his friend Sir Walter Raleigh figure out the best way to stack cannonballs on ship
decks.  The goal was to find the densest possible spherical arrangements~\dots{}  basically,
the way grocers stack oranges''

``Okay,'' he said, nodding.

``Kepler's conjecture {\rm{seems}} perfectly sound,'' I said.

``That is does,'' Ben said.

``But here's the thing\dots.   I looked it up and discovered that, in 1998, a proof
had finally been put forward by an American mathematician named Thomas Hales.  In 2003,
a committee that had been assigned to verify Hales's work confirmed that they were ninety-nine
percent certain of the proof's correctness.  But that one percent was key.  The mathematical
world is still waiting for the publication of the data that will prove the Kepler Conjecture definitively.''

``Sucks for Thomas Hales,'' Ben said.

``I agree.  But it makes sense that they have to be certain, doesn't it?''


{\hfill\rm--Michelle Richmond, No One You Know~\cite{Richmond}} % page 177.

}

}

%%DD Figure of cannonballs.

\bigskip

{

\narrower\parindent=0pt
\parskip=0.4\baselineskip

{\it

``Sometimes fixing a 1 percent defect takes 500 percent effort.''

{\hfill\rm-- Joel Spolsky, Joel on Software~\cite{Spolsky}} % page 122

}

}

\bigskip

{

\narrower\parindent=0pt
\parskip=0.4\baselineskip

{\it

``Every one fully persuaded is a fool.''

{\hfill\rm-- Barthasar Graci\'an, the Art of Worldly Wisdom~\cite{Gracian}} % p110

}

{

\it

%Believe not everything, but only what is proven: the former is foolish, 
%the latter the act of a sensible man. 

%{\hfill -- Democritus}

}

}



\newpage

\section*{The Kepler Conjecture}

In 1611, Johannes Kepler wrote a booklet in which he asserted that the familiar
cannonball arrangement of congruent balls in space achieves the
highest possible density.  This assertion has become known as the
Kepler conjecture.  This book presents a proof.

As early as 1831, Gauss established a special case of the conjecture,
by proving that the cannonball arrangement is optimal among all
\newterm{lattices}~\cite{Gau31}.  Later in the nineteenth century,
Thue solved the corresponding problem in two dimensions, showing that
the hexagonal arrangement of disks in the plane achieves optimal
density~\cite{Thu92}, \cite{Thu10}.  Hilbert, in his famous list of
mathematical problems, made the Kepler conjecture part of his eighteenth
problem.  In 1953, Fejes T\'oth formulated a general strategy to
confirm the Kepler conjecture, but lacked the computational resources
to carry it out~\cite{Fej53}.  The conjecture was finally resolved in
1998, even though the full proof was not published until
2006~\cite{Hales:2006:DCG}.  Section~\ref{sec:history} gives
additional historical background.

The Kepler conjecture has become a test of the capability of computers to
deliver a reliable mathematical proof.  The original proof 
involved many long computer calculations that
exhausted the efforts of a team of referees.  This book has
 redesigned the proof in a way that makes the correctness of
the computer proof as transparent as possible.


\section*{Formal Proofs}

After all is said and done, a proof is only as reliable as the
processes that are used to verify its correctness.  The ultimate
standard of proof is a formal proof, which is nothing other
than an unbroken chain of logical inferences from an explicit set of
axioms.  While this may be the mathematical ideal of proof, actual
mathematical practice generally deviates significantly from the ideal.



%More than ten years have passed since a proof was first
%obtained. Why give a new presentation of the proof?
%
%The original proof was not widely understood.  The complexity was not
%because of conceptual challenges.  In fact, the proof makes only
%modest demands on the theoretical training of the reader.  It is
%possible to read and understand the proof with a knowledge of a
%limited body of mathematics, such as basic calculus and elementary
%Euclidean geometry.
%
%Nevertheless, the proof involves many long calculations. Even worse,
%it it relies on computer calculations.  An error in any calculation or
%a bug in the computer code has the potential to topple the entire
%proof.
%
%The referees were conscientious and checked many of the calculations.
%However, for the most part, the computer code lay beyond the scope of
%referee review, and even careful quality control can let a
%occasional bug slip through undetected.
%
%After all is said and done, no proof is more reliable than the
%reliability of the processes that are used to verify its
%correctness.  These processes include the checking that the author
%makes before releasing the proof for public scrutiny, the checking
%of the referees, and the checking done by readers after publication.

In recent years, as part of this project, I have been increasingly preoccupied by the
processes that mathematicians rely on to ensure the correctness of complex
proofs. A century ago, Russell's paradox and other antinomies threatened
 set theory  with fires of destruction.
Researchers from Frege to G\"odel solved the problem of
rigor in mathematics and found a theoretical solution but did not
extinguish the fire at the foundations of mathematics
because they omitted the practical implementation. Some, such as
Bourbaki, have even gone so far as to claim that ``formalized
mathematics cannot in practice be written down in full'' and call
such a project
``absolutely unrealizable'' \cite[pp.~10--11]{Bour:68:Sets}. % Theory of
                                                          % Sets, page
                                                          % 10,11.

While it is true that formal proofs may be too long to print,
computers -- which do not have the same limitations as paper -- have
become the natural host of formal mathematics. In recent decades,
logicians and computer scientists have reworked the foundations of
mathematics, putting them in an efficient form designed for real use
on real computers.

For the first time in history, it is possible to generate and verify
every single logical inference of major mathematical theorems.  This
has now been done for many theorems, including the four-color theorem, the prime number
theorem, the Jordan curve theorem, the Brouwer fixed point theorem,
and the fundamental theorem of calculus.  Freek Wiedijk
reports that 83\% of a list of one hundred famous theorems have now been
checked formally \cite{wiedijk:100}.  The list of seventeen remaining
theorems contains two particular challenges: the independence of the
Continuum Hypothesis and Fermat's Last theorem.
% list as of 5/5/2010.

Some mathematicians remain skeptical of the process because computers
have been used to generate and verify the logical inferences.
Computers are notoriously imperfect, with flaws ranging from software
bugs to defective chips.  Even if a computer verifies the inferences,
who  verifies the verifier, or then verifies the verifier of the
verifier?  Indeed, it would be unscientific of us to place an
unmerited trust in computers.

The choice comes down to two competing verification processes.  The
first is the traditional process of referees, which depends largely on
the luck of the draw -- some referees are meticulous, while others are
careless.  The second process is formal computer verification, which
is less dependent on the whims of a particular referee.  In my view,
the choice between the conventional  process by human referee and computer
verification is as evident as the choice between a sundial and an atomic
clock in contemporary science.

The standard of proof I have adopted is the highest scientific
standard available by current technology.  
%That standard, which
%continues to evolve with the advancement of technology, is formal
%verification by computer.
%The boundary that separates an easy proof from a difficult
%one shifts with current technology.  
The introduction of steel in
architecture is not a mere reinforcement of wood and stone; it changes
the world of structural possibilities.  There is no longer any
reason to limit proofs to ten-thousand-pages when our
technology supports a million pages.

The style of formal proofs is different from that of conventional
ones.  It is easier to formalize several short snappy proofs
than a few intricate ones.  Humans enjoy surprising new perspectives,
but computers benefit from repetition and standardization.  Despite
these differences, I have sought  proofs that might bring
pleasure to the human reader while providing precise instructions for
the implementation in silicon.

\section*{Conventions}

To make formalization proceed more smoothly, long proofs have been
broken into a series of smaller claims.  Each claim starts a new paragraph
and is set in italics.  The second sentence of the paragraph begins with
the word \newterm{indeed} when the proof of the claim is direct and with
the word \newterm{otherwise} when the proof is indirect by contradiction.

The pronoun \newterm{we} is used inclusively for the author and reader as
we work our way through the proofs in this book.  The pronoun \newterm{I} 
refers to the author alone, as in the first sentence of the following section.

The asterisk $\wild$ is used as a wildcard symbol in patterns.  It
replaces a term in contexts where the name of the term is not
relevant.  It can also denote a bound variable.  For example, the
function $f(\wild,y)$ of a single variable is obtained from $f$ by
evaluating the second argument at a fixed value $y$.

The union of the family $X$ of sets is written as $\bigcup\ X$ or as
$\bigcup_{x\in X} x$ without any difference in meaning.  The first form
is preferred because of its economy.  We also use both expressions
$\bigcap\ X$ and $\bigcap_{x\in X} x$ for the intersection of a family of sets.

The documentation of the computer calculations for the Kepler
conjecture has evolved over time.  The 1998 preprint version of the
proof of the Kepler conjecture contains long appendices of doubtful
value that list hundreds of calculations that enter into the proof.
These appendices were cut from the published version of the proof
because it is more useful to store the computer part of the proof at a
computer code repository that is permanent, versioned, and freely
available.  The computer code and documentation are housed at
\newterm{Google Code project hosting}.  A separate guide, which is
available at the project site, describes the computer calculations
that appear in this book.  When this book uses an external
calculation, it is marked in italic font as a \cc{This explains notation}.

\section*{A Blueprint}

I dream of a fully formally verified solution to the packing problem.
This project is still unfinished, but significant progress is being
made.  The book is a blueprint for formal proofs because it gives the
design of the formal proof to be constructed.  Decisions about what to
include in this book have been shaped by the list of theorems already
available in the library of the proof assistant \newterm{HOL Light}.
For example, this book accepts basic point-set topology and measure
theory because they have been formalized by Harrison~\cite{HOLL}.


The book is divided into four parts, the first of which
 describes the major ideas, methods, and
organization of the proof.  

%There is an essay on each major computer
%component of the proof. The purpose is to provide a panoramic view of
%proof, to provide intuition about proof strategies.  After reading
%this part of the book, the reading should understand what the proof is
%all about, without yet dipping into technical details.

The part on foundations provides background material about
constructions in discrete geometry.  The first of these chapters
covers trigonometric identities and basic vector geometry.  The second
treats volume from an elementary point of view.  The third chapter
covers planar graph theory from a purely combinatorial perspective.
The fourth chapter continues with planar graphs, now from a geometric
perspective.

The next part of the book gives the solution to the packing problem.
The first chapter  gives a top-level overview of the major
steps of the proof,   describing how the problem can be reduced from
a problem with infinitely many variables to one in finitely many
variables.  The remaining chapters in this part flesh out the proof.

The final part of the book views dense sphere packings from a larger perspective.
One chapter resolves some other longstanding conjectures in
discrete geometry: Bezdek's strong dodecahedral conjecture and Fejes
T\'oth's full contact conjecture.  Another chapter describes  formal proof, 
the HOL Light proof assistant, and the Flyspeck formalization project.
An appendix  gives an introduction  to the methods used in the computer calculations.

\section*{Simplifications}

Many simplifications of the original proof have been found over the
past several years.  These simplifications are published here for the
first time.  Gonthier has reworked the proof of the four-color
theorem to avoid the use of the Jordan curve theorem, using instead
the much simpler notion of M\"obius contour from the theory of
hypermaps.  I have followed Gonthier's lead.

The optimality of the face-centered cubic packing is an assertion
about infinite space-filling packings.  For computational purposes, it is
useful to reduce the sphere packing problem to finite packings.  A
\newterm{correction term} is associated with each different reduction from
infinite packings to finite packings.  Ferguson and I worked together to
find the original proof of the Kepler conjecture.  The two of us considered a
large number of different correction terms, seeking one that
would simplify the computations as much as possible.  In a discussion
of the solution of the packing problem, I wrote that ``correction
terms are extremely flexible and easy to construct, and soon Samuel
Ferguson and I realized that every time we encountered difficulties in
solving the minimization problem, we could adjust $f$ [the correction
term] to skirt the difficulty\dots.  If I were to revise the proof
to produce a simpler one, the first thing I would do would be to
change the correction term once again.  It is the key to a simpler
proof''~\cite{Hales:2000:cannonballs}.  Marchal has recently found a remarkably simple 
correction term, giving a new way  to  reduce from infinite packings
to finite packings~\cite{marchal:2009}.  This book implements his reduction step.

There are many other improvements of the proof that are not visible in
the book because they are implemented in computer code, including
a reduction of  the number of lines of computer code from over 187,000
to well under 10,000.  Needless to say, the quickest way to verify that a block
 of computer code is free of bugs is to delete the code altogether.




\bigskip
\hbox{}



\bigskip
\hbox{}

{
\parindent=0pt
\obeylines

Thomas C. Hales
Pittsburgh, PA
May 2010

}








 
%%%%%%%%%%%%%%%%%%%%%%%%%%%%%%%%%%%%%%%%%
  \mainmatter
    


\def\odpcvgh{
\tikzfig{trig}{\guid{ODPCVGH} Trigonometric and inverse trigonometric
functions}
%
%arctangent function on the domain \leftopen -4,4\rightopen\ 
%and the $\arccos$ function on $\leftclosed-1,1\rightclosed$.}}
{
[scale=0.5]
\draw (-2*1.57,0) sin (-1.57,-1) cos (0,0) sin (1.57,1) cos (3.14,0) sin (3*1.57,-1);
\draw   (-2*1.57,-1) cos (-1.57,0) sin  (0,1) cos (1.57,0) sin (2*1.57,-1) cos (3*1.57,0); 
\draw[help lines,<->] (-3.3,0) -- (3*1.57 + 0.2,0);
\draw[help lines,<->] (0,-1) -- (0,2.0);
\draw plot[smooth] file {tikz/tan.table};
\node at (-0.5,-1.8) {$\tan$};
\node at (2,0.5) {$\sin$};
\node at (1.1,-0.3) {$\cos$};
% GG need axis labels and ticks, base points of labels should be precisely aligned.
\begin{scope}[xshift=10cm]
\draw plot[smooth] file {tikz/arctan.table} node[above] {$\arctan$};
\draw plot[smooth] file {tikz/arccos.table} node[right] {$\arccos$};
%\draw[gray,->,very thin] (-1.2,0) -- (1.4,0);
\draw[help lines,<->] (0,-1.6) -- (0,3.2);
\draw[help lines,<->] (-4,0) -- (4,0);
\end{scope}
}}

%
\tikzfig{circle}{\guid{HGMTQFG} Lemma~\ref{lemma:circle} as a special case of the Pythagorean theorem}
{
[scale=0.1]
\draw (0,0)  --(12,0) --  (12,5) --  cycle;
\draw[very thin] (11,0) -- (11,1) -- (12.0,1);
\path (5,-1.5) node {$\cos x$};
\path (16,2.5) node {$\sin x$};
\path (6,5)  node {$1$};
}

%
\tikzfig{tan}{\guid{GQQAKYI} 
The tangent function on $\leftopen-\pi/2,\pi/2\rightopen$}
{
[scale=0.2]
\draw plot[smooth] file {tikz/tan.table};
\draw[help lines,->] (-1.57,0) -- (1.57,0);
\draw[help lines,->] (0,-6.0) -- (0,6.0);
}

\tikzfig{arctrig}{\guid{RUJPPWJ} 
The arctangent function on the domain \leftopen -4,4\rightopen\ 
and the $\arccos$ function on $\leftclosed-1,1\rightclosed$.}
{
[scale=0.4]
\draw plot[smooth] file {tikz/arctan.table};
\draw plot[smooth] file {tikz/arccos.table};
%\draw[gray,->,very thin] (-1.2,0) -- (1.4,0);
\draw[help lines,->] (0,-0.2) -- (0,3.2);
\draw[help lines,->] (-4,0) -- (4,0);
}

\tikzfig{atn-polar}{\guid{YOXQFUB} 
The function $\atn$ gives the polar angle $\theta$ of $(x,y)$.}
{
[scale=0.15]
\draw[gray,->,very thin] (-4,0) -- (14,0);
\draw[gray,->,very thin] (0,-2) -- (0,5);
\draw (0,0)  --(12,0) --  (12,5) --  cycle;
\draw[very thin] (11,0) -- (11,1) -- (12.0,1);
\path (6,-1.5) node {$x$};
\draw[very thin] (4,0) arc (0:22.62:4);
\path (14,2.5) node {$y$};
\path (2,3) node {$\theta$};
}

\def\odpcvgh{
\tikzfig{trig}{\guid{ODPCVGH} Trigonometric and inverse trigonometric
functions}
{
[scale=0.5]
\draw[help lines,<->] (-3.3,0) -- (3*1.57 + 0.2,0);
\draw[help lines,<->] (0,-1) -- (0,2.0);
\draw plot[smooth] file {tikz/tan.table};
\node at (-0.5,-1.8) {$\tan$};
\node at (2,0.5) {$\sin$};
\node at (1.1,-0.3) {$\cos$};
% GG need axis labels and ticks, base points of labels should be precisely aligned.
\begin{scope}[xshift=10cm]
\draw plot[smooth] file {tikz/arctan.table} node[above] {$\arctan$};
\draw plot[smooth] file {tikz/arccos.table} node[right] {$\arccos$};
%\draw[gray,->,very thin] (-1.2,0) -- (1.4,0);
\draw[help lines,<->] (0,-1.6) -- (0,3.2);
\draw[help lines,<->] (-4,0) -- (4,0);
\end{scope}
}}

%
\tikzfig{circle}{\guid{HGMTQFG} Lemma~\ref{lemma:circle} as a special case of the Pythagorean theorem}
{
[scale=0.1]
\draw (0,0)  --(12,0) --  (12,5) --  cycle;
\draw[very thin] (11,0) -- (11,1) -- (12.0,1);
\path (5,-1.5) node {$\cos x$};
\path (16,2.5) node {$\sin x$};
\path (6,5)  node {$1$};
}

%
\tikzfig{tan}{\guid{GQQAKYI} 
The tangent function on $\leftopen-\pi/2,\pi/2\rightopen$}
{
[scale=0.2]
\draw plot[smooth] file {tikz/tan.table};
\draw[help lines,->] (-1.57,0) -- (1.57,0);
\draw[help lines,->] (0,-6.0) -- (0,6.0);
}

\tikzfig{arctrig}{\guid{RUJPPWJ} 
The arctangent function on the domain \leftopen -4,4\rightopen\ 
and the $\arccos$ function on $\leftclosed-1,1\rightclosed$.}
{
[scale=0.4]
\draw plot[smooth] file {tikz/arctan.table};
\draw plot[smooth] file {tikz/arccos.table};
%\draw[gray,->,very thin] (-1.2,0) -- (1.4,0);
\draw[help lines,->] (0,-0.2) -- (0,3.2);
\draw[help lines,->] (-4,0) -- (4,0);
}

\tikzfig{atn-polar}{\guid{YOXQFUB} 
The function $\atn$ gives the polar angle $\theta$ of $(x,y)$.}
{
[scale=0.15]
\draw[gray,->,very thin] (-4,0) -- (14,0);
\draw[gray,->,very thin] (0,-2) -- (0,5);
\draw (0,0)  --(12,0) --  (12,5) --  cycle;
\draw[very thin] (11,0) -- (11,1) -- (12.0,1);
\path (6,-1.5) node {$x$};
\draw[very thin] (4,0) arc (0:22.62:4);
\path (14,2.5) node {$y$};
\path (2,3) node {$\theta$};
}




  % restore this
    %%\chapter{Essays}


%% XX Write essay on Formal Proof.
%% XX Finish Mathematica calc on dodec.

\chapter{Close Packing}
\section{Face Centered Cubic}



The face-centered cubic packing is the familiar pyramid arrangement
of balls on a square base.  It is also the pyramid arrangement on 
a triangular base.  The square base packing and triangular base
packing differ only in their orientation in space.  
Figure~\ref{fig:tri-square-cannonballs} 
shows how the triangular base packing fits between the peaks
of two adjacent square base pyramids.

%% WW not yet done.
\begin{figure}[htb]
  \centering
  \myincludegraphics{noimage.eps}
  \caption{The square and triangular based packings of cannonballs}
  \label{fig:tri-square-cannonballs}
\end{figure}

Density, defined as a ratio of volumes, is insensitive to changes
of scale.  For convenience, we will always work with balls of unit radius.
This means that the distance between centers of balls in a packing
will always be at least $2$.  We identify a packing with
its set of centers.  Thus, for our purposes a packing is just
a set of points in $\ring{R}^3$, whose elements are separated by
distances at least $2$.


% based on page 28 of kepler research scan/scan0001.tif
The fcc packing is an alternating tiling of regular
tetrahedra and regular octahedra.  See Figure~\ref{fig:tet-oct}.
%If the balls have
%unit radius, each tetrahedron and octahedron has edge length $2$.
The figure has one
tetrahedron at each vertex and one octahedron in the center.
By similarity, the total volume is $8 = 2^3$ times the volume of each
smaller tetrahedron.
This dissection exhibits
the volume of a regular octahedron as exactly four times
the volume of a regular tetrahedron of the same edge length.

%% WW not yet done.
\begin{figure}[htb]
  \centering
  \myincludegraphics{noimage.eps}
  \caption{The fcc as a tiling of tetrahedra and octahedra}
  \label{fig:tet-oct}
\end{figure}









The face-centered cubic packing is the packing obtained
from a cubic lattice, with a ball at each vertex and adding
a ball at the center of each face.  The name `face-centered cubic' comes from this construction.  The edge of each cube is
$2\sqrt2$, and the diagonal of each face $4$, for balls
of unit diameter.  The density of the packing is equal to the density
within each cube.  The cube has volume $(2\sqrt2)^3 = 16\sqrt2$.
It contains a total of four balls:  half a ball along each
of six faces and one eighth a ball at each of eight corners.
Thus, the density is
   $$
   \frac{   4 (4\pi/3)}{16\sqrt2} = \frac{\pi}{\sqrt{18}}.
   $$
See Figure~\ref{fig:fcc-cube}.

%% WW not yet done.
\begin{figure}[htb]
  \centering
  \myincludegraphics{noimage.eps}
  \caption{The face-centered cubic structure}
  \label{fig:fcc-cube}
\end{figure}


The alternating arrangement of regular tetrahedra and octahedra can be
superimposed on the cubic picture.  There is one tetrahedron at each vertex
of the cube, extending to the centers of the three adjacent faces.
An octahedron sits the center of the cube, a vertex
at the center of each face.  There is an additional quarter of an
octahedron along each edge, extending to the midpoints of the two adjacent
faces.  This gives a total of eight tetrahedra and four octahedra.  
As each
octahedron has the volume of four tetrahedra, exactly $1/3$ of
the cube is filled with tetrahedra, the other $2/3$ with octahedra.
This decomposition determines the volume $2\sqrt2/3$ of a tetrahedron.
%(pretend we didn't know). The volume $16\sqrt2$ 
%of the cube equals $24$ tetrahedra $\ldots$, giving each a volume
%of 
%$2\sqrt{2}/3$.

The density of the face-centered cubic packing is the weighted density
of the densities of the tetrahedron and octahedron.  Write $\dtet$
and $\doct$ for these densities.  For example, $\dtet$ is the ratio
of the volume
of the part within the tetrahedron of the unit balls (at the four vertices)
to the volume of the tetrahedron.  As the cube is $1/3$ filled  with
tetrahedra and $2/3$ filled with octahedra, we have
$$
  \frac{\pi}{\sqrt{18}} = \frac{1}{3}\dtet + \frac{2}{3}\doct.
$$

The Voronoi cell of a point $v$ 
in a packing $\Lambda$ consists of
all points in $\ring{R}^3$ that are closer to $v$ than to any
other point of $\Lambda$.
Each Voronoi cell of the face-centered cubic
packing is a rhombic dodecahedron (Figure~\ref{fig:rhombic}).   %% not drawn.
The rhombic dodecahedron can be constructed from a cube by placing a
square based pyramid (with height half as great as an edge of its square
base) on each of the six faces (Figure~\ref{fig:rhombic-cube}).  %% not drawn. 
For the scale to be correct, the
initial cube should have edge $\sqrt{2}$. 

The Voronoi cells can be superimposed on the cube of side $2\sqrt2$.  If we center
a cube of side $\sqrt2$ at each ball in the packing. These cubes fill
the black grid of an infinite three dimensional checkerboard,
leaving a second grid of white cubes unfilled.  Each
white cube can be partitioned into  pyramids along
its faces with apex at the center of the cube.  Attaching these pyramids
along their bases to the adjoining black cubes gives the Voronoi cell.

%% WW not yet done.
\begin{figure}[htb]
  \centering
  \myincludegraphics{noimage.eps}
  \caption{The rhombic dodecahedron as Voronoi cell}
  \label{fig:rhombic}
\end{figure}

%% WW not yet done.
\begin{figure}[htb]
  \centering
  \myincludegraphics{noimage.eps}
  \caption{The rhombic dodecahedron constructed from a cube}
  \label{fig:rhombic-cube}
\end{figure}

Each Voronoi cell contains one black cube of side $\sqrt2$ and a total
of one white cube, for a total volume of $4\sqrt2$.  This constant is one of the
fundamental constants in this book.  The volume of every Voronoi cell
is compared against the volume of the rhombic dodecahedron.
The density of the face-centered
cubic is the ratio of the volume of a ball to the volume of its Voronoi
cell, which gives $\pi/\sqrt{18}$, yet again.





\section{Hexagonal-close packing}

There
is a popular and persistent misconception in the popular press
that the face-centered cubic packing is the only packing with
density $\pi/\sqrt{18}$. 

In the face-centered cubic packing, each ball is tangent to twelve
others.  For each ball in the packing, this arrangement of twelve
tangent balls is the same.  We call it the fcc pattern. 
In the
hexagonal-close packing, each ball is tangent to twelve others
(Figure~\ref{fig:hcp}).
For each ball in the packing, the arrangement of twelve tangent
balls is again the same.  We call it the hcp pattern.  The fcc
pattern is different from the hcp pattern.  In the fcc pattern,
there are four different planes through the center of the central
ball that contain the centers of six other balls at the vertices
of a regular hexagon.  In the hcp pattern, there is only one such
plane.  We call the arrangement of balls tangent to a given ball
the {\it local tangent arrangement} of the ball.

%% WW not yet done.
\begin{figure}[htb]
  \centering
  \myincludegraphics{noimage.eps}
  \caption{The fcc and hcp patterns}
  \label{fig:hcp}
\end{figure}

There are uncountably many packings of density $\pi/\sqrt{18}$
that have the property that every ball is tangent to twelve others
and such that the tangent arrangement around each ball is either
the fcc pattern or the hcp pattern.

By {\it hexagonal layer}, we mean a translate of the two-dimensional
lattice of points $M$ in the $A_2$ arrangement. That is, $M$ is a
translate of the planar lattice generated by two vectors of length
$2$ and angle $2\pi/3$.  The face-centered cubic packing is an
example of a packing built from hexagonal layers.

If $M$ is a hexagonal layer, a second hexagonal layer $M'$ can be
placed parallel to the first so that each lattice point of $M'$ has
distance $2$ from three different vertices of $M$.  When the second
layer is placed in the manner, it is as close to the first layer as
possible. Fix $M$ and a unit normal to the plane of $M$. The normal
allows us to speak of the second layer $M'$ as being ``above'' or
``below'' the layer $M$. There are two different positions in which
$M'$ can be placed closely above $M$ and two different positions in
which $M'$ can be placed closely below $M$  (Figure~\ref{fig:two-holes}). As we build a packing,
layer by layer, ($M$, $M'$, $M''$, and so forth), there are two
choices at each stage of the close placement of the layer above the
previous layer. Running through different sequences of choices gives
uncountably many packings.  In each of these packings the tangent
arrangement around each ball is that of the twelve balls in the
face-centered cubic or the twelve balls in the hexagonal-close
packing.

%% WW not yet done.
\begin{figure}[htb]
  \centering
  \myincludegraphics{noimage.eps}
  \caption{The two positions of hexagonal layers}
  \label{fig:two-holes}
\end{figure}

Let $\Lambda$ be a packing built as a sequence of close-packed
hexagonal layers in this fashion.  If $P$ is any plane parallel to
the hexagonal layers, then there are at most three different
orthogonal projections of the layers $M$ to $P$.  Call these
projections $A$, $B$, $C$.  Each hexagonal layer has a different
projection than the layers immediately above and below it.  In the
fcc packing, the successive layers are $A,B,C,A,B,C,\ldots$.  In
the hcp packing, the successive layers are $A,B,A,B,\ldots$.  If
we represent $A$, $B$, and $C$ as the vertices of a triangle, then
the succession of hexagonal layers can be described by a walk
along the vertices of the triangle. Different walks through the
triangle describe different packings (Figure~\ref{fig:fcc-tri}).

%% WW not yet done.
\begin{figure}[htb]
  \centering
  \myincludegraphics{noimage.eps}
  \caption{Describing layers by walks through a triangle}
  \label{fig:fcc-tri}
\end{figure}

% In the face-centered
%cubic packing, every local tangent arrangement of $12$ tangent
%balls is the same.  We call this local tangent arrangement the
%fcc-pattern. Similarly, the local tangent arrangement of the
%hexagonal-close packing will be called the hcp-pattern.  These
%patterns are uniquely determined up to rigid motion.

The different walks through a triangle give all packings of
infinitely many equal balls in which the tangent arrangement around
every ball is either the fcc pattern of twelve balls or the hcp
pattern of twelve balls.



We prove that different walks through a triangle give all
such packings. Assume first 
that $\Lambda$ does not contain any balls whose local
tangent arrangement is the hcp pattern.  Then every local tangent
arrangement is the fcc pattern, and $\Lambda$ itself is then the
face-centered cubic packing.  This completes the proof in this case.

Now assume that a packing $\Lambda$ contains some ball
(centered at $v_0$) in the hcp pattern. The hcp pattern contains a
uniquely determined plane of symmetry. This plane contains $v_0$ and
the centers of six others arranged in a regular hexagonal. If $v$ is
the center of one of the six others in the plane of symmetry, its
local tangent arrangement of twelve balls must include $v_0$ and an
additional four of the twelve balls around $v_0$. These five centers
around $v$ are not a subset of the fcc pattern. They can be uniquely
extended to twelve centers arranged in the hcp pattern. This hcp
pattern has the same plane of symmetry as the hcp pattern around
$v_0$. In this way, as soon as there is a single center with the hcp
pattern, the pattern propagates along the plane of symmetry to
create a hexagonal layer $M$.

Once a packing $\Lambda$ contains a single hexagonal layer, the
condition that each ball be tangent to twelve others forces a
hexagonal layer $M'$ above $M$ and another hexagonal layer below
$M$.  Thus, a single hexagonal layer forces an infinite sequence of
close-packed hexagonal layers.  This completes the proof.



\section{Gauss}

Gauss proved that the face-centered packing has the greatest density
of any lattice packing in three dimensional Eulidan space.  
There is a short proof that
does not require any calculations.

Start with an arbitrary lattice $\Lambda$ in which every point has
distance at least two from every other.  Center a unit ball at each
point in the lattice.  A lattice of greatest density has the property that some pair of balls touch.  The lattice property
then forces the the balls into parallel infinite linear strings
like beads on a string.  The lattice of greatest
density has the property that two of these 
infinite parallel strings touch.  The lattice property then arranges the strings into parallel sheets.  On each sheet the touching parallel
strings form a rhombic tiling.  The lattice of greatest density
has the property that each parallel sheet should sit
as snugly as possible on the sheet below.  That means that some ball (centered at $A$) of
one sheet touches three balls (centered at $B,C,D$) 
on a the next layer down (Figure~\ref{fig:rhombus}).

%% WW not yet done.
\begin{figure}[htb]
  \centering
  \myincludegraphics{noimage.eps}
  \caption{A ball from one sheet touches three balls centered at points of a rhombus.}
  \label{fig:rhombus}
\end{figure}

As the balls on each sheet form a rhombic tile, two of the distances
between $B,C,D$, corresponding to two edges of the rhombus, 
are equal to $2$.  This means that $A$ together with
two of $B,C,D$ form an equilateral triangle.  From the perspective of the plane
containing this equilateral triangle, the lattice property forces this entire
plane, as well as parallel planes, to be tiled with equilateral triangles.
From the earlier argument, each plane sits as snugly as possible on
the sheet below.  Some ball of one sheet touches the three balls forming
an equilateral triangle in the layer below.  These four balls form a regular
tetrahedron.  This tetrahedron 
uniquely identifies the lattice as the face-centered cubic.






\section{Thue}\label{sec:thue}


As mentioned, 
Thue solved the  packing problem for congruent disks in the plane.
The optimal packing is the hexagonal packing (Figure~\ref{fig:2D-hex}). 
The density of this packing is $\pi/\sqrt{12}$: the ratio of the
area of a unit disk to the area of a hexagon of inradius one.
It admits an elementary solution that we sketch here.  B. Casselman
has an animated interactive demo of this solution \cite{BC00}.

%% WW not yet done.
\begin{figure}[htb]
  \centering
  \myincludegraphics{noimage.eps}
  \caption{The optimal packing in two dimensions}
  \label{fig:2D-hex}
\end{figure}

Let $\Lambda$ be the set of centers of a collection of unit disks.
Take the Voronoi cell around each disk.  It is enough to show that
each Voronoi cell has area at least $\sqrt{12}$.  For simplicity,
we may assume that $0\in\Lambda$ is the center of our Voronoi cell.  

We truncate the Voronoi cell by intersecting it with a disk of radius
$\rho=2/\sqrt3$.  It is enough to show that the truncated Voronoi cell has
density  at most $\pi/\sqrt{12}$.  

There is not a point $w$ in the plane that has distance less than $\rho$
from three disk centers $v_1,v_2,v_3$.  Otherwise, one of the three
angles $\gamma$ at $w$ is at most $2\pi/3$; and $\cos\gamma\ge -0.5$  
The law of cosines applied
to the triangle $w,v_i,v_j$ with angle $\gamma$ gives the contradiction:
   $$
   4 \le c^2 = a^2 + b^2 - 2 a b \cos\gamma 
   \le a^2 + b^2 + a b < 3\rho^2 = 4.
   $$
Thus, the boundary of the truncated Voronoi cell consists of circular
arcs and chords of the circle of radius $\rho$ (see Figure~\ref{fig:2D-proof}).

%% WW not yet done.
\begin{figure}[htb]
  \centering
  \myincludegraphics{noimage.eps}
  \caption{The sectors and triangles of a truncated Voronoi cell}
  \label{fig:2D-proof}
\end{figure}

The
parts of the Voronoi cell that lie within a circular sector have density
$1/\rho^2 = 3/4 < \pi/\sqrt{12}$.  The
parts of the Voronoi cell that
lie within a triangle have density
   \begin{equation}\label{eqn:rog2d}
   \frac{\theta}{\rho^2 \cos\theta\sin\theta}
   \end{equation}
% checked 3/31/2008.
where $0 \le \theta\le \pi/6$.  An easy optimization gives the maximum
at $\theta=\pi/6$ with value $\pi/\sqrt{12}$.

In some ways it it unfortunate that the problem in two dimensions is so
elementary.  It gives meager hints about how to solve the problem in
three dimensions.   It suggests the
introduction of Voronoi cells and the usefulness of truncation.
The little optimization problem on triangles in Equation~\ref{eqn:rog2d}
generalizes to $n$-dimensions.  This is Rogers's lemma.
It appears as Lemma~\ref{lemma:rogers} in this book.  It is a tremendously useful lemma in the study of packing densities.

%\subsection{two dimensions}

\bigskip

There are many proofs of Thue's theorem.  Here we present a second proof
that is due to L. Fejes T\'oth.  Section~\ref{sec:13s} 
introduces the Delaunay
triangulation of a finite set $C$ of points on a sphere.
There is a corresponding Delaunay triangulation for infinite sets $\Lambda$
of points in the plane (or more generally in $n$-dimensions).  We
assume that $\Lambda$ is the set of centers of packing so that
distances between points of $\Lambda$ are at least two.  
If we enlarge a packing
$\Lambda\subset\Lambda'$, then the density can only go up.  In issues
of density, it does no harm to assume that the packing is {\it saturated}; that is, it is not a proper subset of another packing.  A saturated
packing of the plane (or $n$-space) has a Delaunay triangulation.  It
is a triangulation in which the circumscribing circle of each triangle
contains no points of $\Lambda$ in its interior.   Each Delaunay
triangle has
radius at most $2$, for otherwise the packing is not saturated: an additional point can be placed at the center of the circumscribing circle.

\begin{proof}
Admitting the existence of the Delaunay triangulation, the proof 
of the packing problem in two dimensions is
elementary.  Each Delaunay triangle contains a portion of a disk at each of
its three vertices.  The interior angles of a triangle sum to $\pi$, giving
half a disk per triangle.  If we show that each triangle has area at least
$\sqrt{3}$, then the density of the packing is at most
 $(\pi/2)/(\sqrt{3}) = \pi/\sqrt{12}$.  To minimize the area of Delaunay
triangle $ABC$, we first replace it with a smaller similar triangle whose
shortest edge (say $BC$) has length $2$.  The third vertex $A$ is constrained
to have distance at least $2$ from $B$ and $C$, and to give circumradius at most $2$.  The constraints on $A$ form three circular
arcs as shown in Figure~\ref{fig:2D-FT} % not drawn.

%% WW not yet done.
\begin{figure}[htb]
  \centering
  \myincludegraphics{noimage.eps}
  \caption{The constraints on the point $A$, vertex of a Delaunay triangle}
  \label{fig:2D-FT}
\end{figure}

The minimizing triangle is determined by the point $A$ closest to the
line $BC$.  There are three such points, all giving triangles of area
exactly $\sqrt3$.  This completes the proof.
\end{proof}

This proof suggests that the solution to the packing problem
in three dimensions might also make use of the Delaunay
decomposition.  This was the starting point of my work on the
problem, to find a three dimensional approach to the packing
problem that generalizes the proof just presented \cite{remarks}.



\clearpage







    
%    \def\achapter#1{\chapter{#1}}
%\def\asection#1{\section{#1}}


    
\def\asection#1{\chapter{#1}}
\def\asubsection#1{\section{#1}}

  \begin{abstract} These are unpublished appendices for 
 "Dense Sphere Packings" (DSP).
 These are notes that were useful in the formalization of some
 lemmas in the proof of the Kepler conjecture but were
 considered too technical for inclusion in DSP.
 It also includes material related to some changes that were made to
  the formalization that occurred after the publication of DSP.

  These appendices have not been edited into a form that is suitable
   for publication. They were primarily written for the benefit
  of the members of the Flyspeck team who carried out the formalization.
    \end{abstract}


    \linput{supplement_local}
    \linput{supplement_general}
       \linput{supplement_oxl}
        \linput{supplement_hypermap}

\bibliographystyle{plainnat}
\bibliography{all}


\end{document}
